
\documentclass{article}

\usepackage{amsmath, amssymb, amsthm}
\usepackage{graphicx}
\usepackage{hyperref}

\title{Autograd in NumPy}
\author{Luke Zeng}
\date{}

\begin{document}

\maketitle

\section{Introduction}
In this article, we will implement a simple autograd system in NumPy. Autograd is a technique for automatically computing gradients of functions, which is essential for training machine learning models without manually deriving gradients. We will define a class for tensors that can track operations and compute gradients, as well as some basic operations like addition and matrix multiplication. We will also implement a simple ReLU activation function and demonstrate how to use our autograd system to compute gradients for a simple neural network.

\end{document}

